\documentclass[11pt]{article}
% Shared preamble for all daily exercises
% Update this file to change settings (padding, parindent, etc.) for all exercises

\usepackage[margin=0.75in]{geometry}
\usepackage{amsmath}
\usepackage{amssymb}

% Custom settings
\setlength{\parindent}{0pt}
\setlength{\parskip}{0.2cm}

\title{Daily Exercise}
\author{Dhruman Gupta}
\date{26th March}

\begin{document}

% Common header for daily exercises
% This file can be included in other LaTeX documents

\maketitle

\noindent
\textbf{Course:} CS-3330, Theory of Computation

\vspace{0.2cm}

\noindent
\textbf{Question:}\noindent 
Analyze the time complexity of Breadth-First Search on a graph in the Turing machine model.

\textbf{Answer:}

Let $G = (V,E)$ be the input graph. Suppose $|V| = n$ and $|E| = m$. We analyse the usual BFS algorithm in the Turing machine model.

On a RAM, BFS runs in time $O(n+m)$, since each vertex is visited at most once and each edge is examined at most once.

Now consider a Turing machine. The graph is written on the tape, say as adjacency lists. The main issue is that a Turing machine does not have random access memory. So operations such as checking whether a vertex has been visited, marking it as visited, maintaining the queue, or locating the adjacency list of a vertex may require scanning a tape segment of size $O(n+m)$.

Thus, an operation that takes constant time on a RAM may take $O(n+m)$ time on a Turing machine.

Since BFS performs $O(n+m)$ such operations overall, the total running time is
\[
O((n+m)(n+m)) = O((n+m)^2).
\]

Hence Breadth-First Search runs in polynomial time on a Turing machine, namely $O((n+m)^2)$.

\end{document}
